% 本文件由 LaTeX 排版 Agent 根据有效 translation.json 自动生成。
% author=Councils; work_id=3814; title=特鲁罗会议

\chapter{特鲁罗会议}

% structure: p0001 source=h1 target=omitted reason=page-title-already-rendered-by-parent

% structure: p0002 source=h2 target=section reason=relative-heading-rank; base=h2

\section{第一教规}

万善之序，莫过于使一切言语和行为开始并终结于天主。因此，为使虔诚得以清晰展示，并使以基督为根基的教会不断增长与进步，超越黎巴嫩的香柏；如今我们赖天主恩宠，在我们法令的开端，决志：由天主所拣选的宗徒——他们亲见过圣言且为其仆役——所确立的信德，当毫无创新、不变且不被侵犯地保存。\par

此外，三百一十八位圣洁真福教父——他们在我们的皇帝君士坦丁治下于尼西亚集会，对抗不虔的亚略及其教导的异教神祇多样性，或准确地说，众多神祇——凭借信友的一致认同，向我们启示并宣告了包含于天主性体中的三位之同体，不使这信德隐藏于无知的斗下，而是公开教导信友以同一敬拜朝拜圣父、圣子及圣神，驳斥并吹散等级差异的论调，摧毁并倾覆异端者针对正统信仰所捏造的如沙堆般的幼稚玩物——我们也确认该信德。\par

同样，我们也确认由一百五十位教父——他们在我们的皇帝狄奥多西大帝时期于帝都集会——所确立的信德，接受他们关于圣神、宣告其天主性的决定，驱逐亵渎的馬其頓尼（连同之前所有真理的敌人），因为他竟敢判断那为主者为仆人，且如盗贼般欲分裂不可分割的统一体，以致我们信仰的完美奥迹不存。\par

并与这可憎可厌的真理对抗者一同，我们谴责阿波里纳留——相同不义之徒的司铎——他不虔地喷出：主所取的身体没有灵魂，并由此推论他为我们所完成的救恩是不完美的。\par

此外，由二百位负天主者教父——他们在我们的皇帝狄奥多西（阿尔卡狄乌斯之子）时代于厄弗所城所确立的道理，我们赞同为虔诚的不中断力量，教导基督、降生成人的天主子，是唯一的一位；并宣告那无人类种子而生他的，是无玷的终身童贞女，光荣她为真实确切的天主之母。我们谴责聂斯多略的荒谬分裂为与天主计划相悖，他教导那位惟一的基督由一个人和一位神性分别组成，并重燃犹太人的不虔敬。\par

此外，我们确认六十三位蒙天主认可的教父——在我们的皇帝玛尔西安时代于大都市迦克墩——依照正统所确立的信德，他们以宏大响亮的声音传至地极：惟一的基督、天主子，具有两个性体，且须在这两个性体受光荣；并自教会圣所逐出尤提赫如同避之不及的黑色瘟疫——他愚蠢空洞地胡言乱语，教导降生成人的伟大奥迹（\OriginalTerm{救恩计划}{\GreekText{οἰκονωμίας}}）仅在思想中完成。并与他一同（谴责）聂斯多略和狄奥斯库若，前者是分裂的辩护者和捍卫者，后者是混淆（两个性体于一位基督中）者；两人皆从他们不虔敬的歧途堕入共同的沉沦和否认天主的深渊。\par

我们也承认为受圣神感动，一百六十五位负天主者教父——他们在我们真福纪念的皇帝犹斯提尼安时代于帝都集会——的虔诚之声，并将之教导后辈；因为他们藉大公会议绝罚并咒逐了莫普苏埃斯提亚的狄奥多（聂斯多略之师）、奥力振、狄第默及埃瓦格留，他们都重新引入虚构的希腊神话，再次带回某些身体和灵魂的轮回，以及错乱的转世（或轮回）成为他们精神的游荡或幻梦，并不虔地侮辱死者的复活。此外（他们谴责）狄奥多雷所写的反对正统信仰和反对真福济利禄《十二章》的文字，以及那封据说由依巴斯所写的信函。\par

我们也同意保持第六次神圣大公会议的信德不受触动——该会议最初在我们真福纪念的皇帝君士坦丁时代于帝都集会，这信德因虔诚的皇帝用自己的印玺批准为后世保障所写的内容而得到更大确认。该会议教导我们应当公开宣认我们的信德：在我们的真天主耶稣基督的降生成人中，有两个本性意志或意愿和两个本性运作；并以公义的判决谴责那些篡改真道、教导人民在主耶稣基督内只有一个意志和一个运作的人，即：法兰的狄奥多、亚历山大的居鲁士、罗马的霍诺留、塞尔吉、皮尔、保禄和伯多禄（他们曾是这座由天主保护的都城的主教）、安提约基雅的主教玛加略、他的门徒斯德望，以及疯狂的波里赫罗尼，从此剥夺他们与我们天主的基督圣体的共融。\par

一言以蔽之，我们决志：信德当坚定不移、毫无玷污地存至世界终结，连同由天主所传授的著作以及所有美化装饰了天主的教会、成为世上之光、拥抱了生命之言者的教导。我们弃绝并绝罚那些为他们所弃绝绝罚的人，视之为真理之敌、疯狂敌视天主者，以及高举不义者。\par

但凡有人不遵守并接纳上述虔诚的教令，不依此教导和宣讲，反而试图反对，则应依照已获认可的圣善教父所颁布的法令，受\OriginalTerm{弃绝}{anathema}，并应被逐出、被剔除，视为与基督徒团体隔绝。因为我们的法令既不添加先前已定义之事，也不删除什么，我们也没有这种权能。\par

% structure: p0013 source=h2 target=section reason=relative-heading-rank; base=h2

\section{第二条}

此神圣会议亦认为，此前由圣善教父所接纳和批准的八十五条教规，以及以圣光荣宗徒之名传于我们的教规，应从此时起坚定不移，为治愈灵魂和医治混乱。在这些教规中，我们被命令接受由克莱孟编写的《圣宗徒宪章》。但从前，因那些背离信仰之人的影响，某些掺假之物被引入，完全违反虔敬，玷污教会，掩盖了神圣法令当前形式的优雅与美丽。因此，我们拒绝这些《宪章》，以便更好地确保最基督徒羊群的建立与安全；绝不接受异端谬误的产物，而应坚守宗徒纯洁完美的教导。我们也认可我们圣善教父所颁布的所有其他神圣教规，即：由聚集在尼西亚的三百一十八位带天主的圣教父，以及在安西拉、在尼奥凯撒利亚、在冈格拉，此外在叙利亚安提约基雅、在弗里吉亚的劳底西亚所颁布的；还有聚集在这位天佑皇城中的一百五十位教父，以及首次在以弗所都城聚集的二百位教父，和在迦克墩的六百三十位圣善教父。同样，还有在萨尔迪卡和在迦太基的教规，以及在此天佑皇城中在主教\Person{内克塔留}和亚历山大里亚总主教\Person{特奥斐禄}下再次聚集的教规。同样，还有前亚历山大里亚大城总主教\Person{狄约尼削}的教规（即教令书），以及亚历山大里亚总主教、殉道者\Person{伯多禄}的教规；尼奥凯撒利亚主教、显灵迹者\Person{额我略}的教规；亚历山大里亚总主教\Person{阿塔纳削}的教规；卡帕多奇亚凯撒利亚总主教\Person{巴西略}的教规；尼撒主教\Person{额我略}的教规；神学家\Person{额我略}的教规；依科尼雍的\Person{安斐罗希}的教规；亚历山大里亚总主教\Person{弟茂德}的教规；同大城亚历山大里亚总主教\Person{特奥斐禄}的教规；同亚历山大里亚总主教\Person{济利禄}的教规；此天佑皇城宗主教\Person{根纳迪}的教规。此外，还有由阿非利加州总主教、殉道者\Person{西彼廉}及其所辖主教会议颁布的教规，此教规仅依照传于上述主教的习俗，在他们地区内保存。任何人不得违反或无视上述教规，或接受其他教规，即某些企图以真理牟利者假托颁布的教规。若有人被证实对上述任何教规进行革新或试图推翻，则应受该教规所定惩罚，并以此矫正其过犯。\par

% structure: p0015 source=h2 target=section reason=relative-heading-rank; base=h2

\section{第三条}

既然我们虔诚的基督徒皇帝曾向此神圣大公会议提出，要求它为圣职人员名单中的人以及那些将神圣事物传授于他人者提供纯洁性，使他们成为无可指责的侍奉者，堪当奉献伟大天主的祭献——祂既是祭品又是大司祭，这祭献是由理智所领悟的；并要求清除他们因非法婚姻而沾染的污秽；而今，至圣罗马教会的人士意图持守完美严格的规则，而此天佑皇城宝座下的人士则持守仁爱与宽仁的规则，两者如此结合，正如我们的教父所做，也如天主的爱所要求，即不让温和流于放纵，也不让严厉流于苛刻；尤其因为无知之过已到达不少人数，我们规定：凡卷入第二次婚姻、直至去年一月十五日（即第四小纪、第六千一百零九年）仍为罪恶奴役、且未决意悔改者，应受教规性的撤职；但若那些卷入此第二次婚姻之紊乱者，在我们的法令颁布前已承认当行之事，并已断绝其罪恶，远离此非法的结合，或其第二次婚姻的妻子已去世，或他们由于学习节制而自愿悔改，迅速忘记先前的不义——无论他们是\OriginalTerm{司铎}{presbyter}或\OriginalTerm{执事}{deacon}——我们决定他们应停止一切司祭侍奉或执行，受一段时间的惩罚，但保留其座位和品位的尊荣，满足于平信徒之前的席位，并含泪恳求上主赦免其无知过犯；因为那自身伤痕待治之人去祝福他人是不适宜的。但那些只娶一妻者，若该妻是寡妇，以及那些在\OriginalTerm{圣秩}{ordination}之后非法进入一次婚姻者，即\OriginalTerm{司铎}{presbyter}、\OriginalTerm{执事}{deacon}和\OriginalTerm{副执事}{subdeacon}，应暂时禁止神圣侍奉并受责罚，然后恢复其原有品位，但不得再晋升任何更高品位，其非法婚姻应公开解除。我们规定此条仅适用于那些在所说的一月十五日（即第四小纪）之前犯有上述过失者；自今时起，我们重新颁布教规，宣告：凡在圣洗后结婚两次、或曾有姘妇者，不能成为主教、\OriginalTerm{司铎}{presbyter}、\OriginalTerm{执事}{deacon}或列入\OriginalTerm{司祭名册}{sacerdotal list}；同样，凡娶寡妇、被休者、娼妓、婢女或伶人者，不能成为主教、\OriginalTerm{司铎}{presbyter}、\OriginalTerm{执事}{deacon}或列入\OriginalTerm{司祭名册}{sacerdotal list}。\par

% structure: p0017 source=h2 target=section reason=relative-heading-rank; base=h2

\section{第四条}

若有任何主教、司铎、执事、副执事、读经员、歌咏员或看门人与献身于天主的女子交合，应予以撤职，因他玷污了基督的净配；但若是平信徒，则予以绝罚。\par

% structure: p0019 source=h2 target=section reason=relative-heading-rank; base=h2

\section{第五条}

凡属神职班者，除教规所列举无可嫌疑之人外，不得拥有任何妇女或女仆，藉以使自己免于任何责难。但若有人违犯此谕令，应予以撤职。阉人亦应遵守同一规则，俾能远见而免受谴责。违犯者，若为圣职人员，予以撤职；若为平信徒，则予以绝罚。\par

% structure: p0021 source=h2 target=section reason=relative-heading-rank; base=h2

\section{第六条}

既然宗徒教规中已声明，在晋升圣职者中，唯有读经员和歌咏员可以结婚；我们亦持守此规定，决定：今后任何副执事、执事或司铎在领受圣秩后，均不得缔结婚姻；若胆敢如此行，应予以撤职。若有任何人进入圣职，愿在领受副执事、执事或司铎之圣职前依法娶妻，则可行。\par

% structure: p0023 source=h2 target=section reason=relative-heading-rank; base=h2

\section{第七条}

我们得知，在某些教会中，执事掌管教会职务，以致有些执事傲慢放肆，竟敢坐在司铎之前。兹决定：即使执事担任显赫职务，即无论身处何种教会职务，均不得在司铎之前就座，除非他代表其宗主教或都主教在另一城市、隶属另一上级时，则因其代表身份应受尊敬。但若有人恃强胆敢如此行，应被逐出其原有等级，在其所属本教会内降为同列最末。因为吾主告诫我们不可贪图首席，此教导见于圣路加福音中，由吾主天主亲自提出。祂对受邀者讲了这比喻：\BibleQuote{凡被人请去赴婚宴，不要坐在首位上，恐怕有比你更尊贵的客被他请来；那请你们的人前来对你说：让座给这一位罢！那时你就要含羞地去坐末位。你被请时，要去坐在末位上，好叫那请你的人来对你说：朋友，请上坐！那时你在同席的人面前就有光彩了。因为凡高举自己的，必被贬抑；凡贬抑自己的，必被高举。}同样规则亦适用于其他神圣等级；因为我们知道属灵之事优于世俗尊荣。\par

% structure: p0025 source=h2 target=section reason=relative-heading-rank; base=h2

\section{第八条}

我们愿我圣教父们所规定的一切事宜得以确立并确认，兹重申该教规：每省主教应在都主教认为最佳之地每年举行会议。但因蛮族入侵及其他偶然原因，教会首长无法每年举行两次会议，故按上述理由，为处理可能发生的教会问题，每省都主教应在其选定之地，于圣复活节至十月之间，每年举行一次上述主教会议。凡主教在各自城市安好、无不可避免之要务而缺席者，应予以兄弟般的责斥。\par

% structure: p0027 source=h2 target=section reason=relative-heading-rank; base=h2

\section{第九条}

不准任何圣职人员开设\Quote{酒馆}。因为连进入酒馆尚且不准，更何况在馆内服侍他人、从事与其身份不相称的生意呢？若有人如此行，应即停止，否则予以撤职。\par

% structure: p0029 source=h2 target=section reason=relative-heading-rank; base=h2

\section{第十条}

主教、司铎或执事若收取利息，即所谓\OriginalTerm{赫卡托斯特}{hecatostæ}者，应立即停止，否则予以撤职。\par

% structure: p0031 source=h2 target=section reason=relative-heading-rank; base=h2

\section{第十一条}

任何圣职人员或平信徒，均不得吃犹太人的无酵饼，不得与其亲密往来，不得在生病时请其治病，不得从其处接受药物，不得与其共浴。若有人胆敢如此行，若为圣职人员，予以撤职；若为平信徒，予以绝罚。\par

% structure: p0033 source=h2 target=section reason=relative-heading-rank; base=h2

\section{第十二条}

此外，我们得知在非洲、利比亚及其他地方，那里最敬爱天主的诸位主教在祝圣后仍与妻子同居，致使百姓跌倒冒犯。因此，我们特别关心一切事物皆有利于托付于我们、交给我们照管的羊群，决定：今后决不容许此类事情发生。我们说这话，并非要废除或推翻古时由宗徒权威所立之事，而是为关心百姓的健康及其向善进益，并免使教会状况受到指责。因为神圣的宗徒说：\BibleQuote{你们或吃或喝，或无论作什么，一切都要为光荣天主而作。你们不可成为使犹太人、或希腊人、或天主的教会跌倒的原因，就如我事事使众人喜欢，不求自己的利益，只求多数人的利益，为使他们得救。你们该效法我，如同我效法基督一样。}若有人被发现如此行，应予以撤职。\par

% structure: p0035 source=h2 target=section reason=relative-heading-rank; base=h2

\section{第十三条}

既然我们知道，罗马教会曾定下规矩，凡堪当升任执事职或司铎职者，须承诺不再与妻子同居；我们则持守古规、宗徒的圆满及秩序，愿自此以后，圣品人的合法婚姻坚定不移，绝不解除他们与妻子的结合，也不在适当的时候剥夺他们彼此的往来。因此，如果有人堪当被祝圣为副执事、执事或司铎，即使他与合法妻子同居，也绝不禁止他获得这样的品级。在祝圣时，也不应要求他承诺戒除与妻子的合法往来，以免我们损害天主所建立、并由祂亲临祝福的婚姻，正如福音所说：\BibleQuote{天主所结合的，人不可拆散；}宗徒也说：\BibleQuote{婚姻应受尊重，床榻应是无玷的；}又说：\BibleQuote{你已与妻子结合吗？不要寻求解脱。}然而我们知道，正如在迦太基集会的人（为圣职人员廉洁的生活而留心）所说的，副执事、执事和司铎，凡接触神圣奥迹的，应按照他们（服务的）次序与配偶禁欲。为此，我们也同样持守宗徒所传承、古习所保存的规矩，知道凡事都有定时，尤其有斋戒和祈祷的定时。因为侍奉天主祭台的，在接触圣物时，应当完全节欲，以便能真诚地向天主求得所求。\par

因此，若有人胆敢违反《宗徒法典》，剥夺任何圣品人——司铎、执事或副执事——与合法妻子的同居与往来，应予以撤职。同样，若任何司铎或执事以虔诚为借口休弃了自己的妻子，应将他逐出共融；若他坚持如此，则应予以撤职。\par

% structure: p0038 source=h2 target=section reason=relative-heading-rank; base=h2

\section{第十四条}

我们神圣的、携带天主的教父的这条法规也应在此得到确认：司铎未满三十岁不得被祝圣，即使他是一位非常值得尊敬的人，也应被暂缓。因为我们的主耶稣基督在三十岁时受洗并开始施教。同样，执事未满二十五岁不得被祝圣，女执事未满四十岁不得被祝圣。\par

% structure: p0040 source=h2 target=section reason=relative-heading-rank; base=h2

\section{第十五条}

\Emph{副执事}未满二十岁不得被祝圣。若有任何人在司祭品级中，违反规定的时间被祝圣，应予以撤职。\par

% structure: p0042 source=h2 target=section reason=relative-heading-rank; base=h2

\section{第十六条}

既然《宗徒大事录》告诉我们，宗徒们选立了七位执事，而新凯撒利亚会议在其公布的法规中，根据《宗徒大事录》规定，即使城市很大，也只应有七位执事；我们征求教父们对宗徒话语的解释，发现他们所说的不是那些在奥迹中服务的人，而是关于伺候饭桌的职务。因为《宗徒大事录》记载如下：\BibleQuote{那时，门徒的人数增多，希腊人向希伯来人发出了怨言，因为他们的寡妇在日常的供应上被忽略了。于是十二位叫齐了众门徒，说：我们放弃天主的圣言而伺候饭桌，实不合宜。所以，弟兄们，从你们当中推选出七位有好名声、且充满圣神和智慧的人，我们派他们管理这一事务。至于我们，我们要专务祈祷，并为圣言服务。这番话使全众都满意，他们就选了充满信德和圣神的斯德望，以及斐理伯、普洛苛洛、尼加诺尔、提孟、帕尔默纳和归依犹太教的安提约基雅人尼苛劳，将他们安置在宗徒面前。}\par

教会圣师金口圣若望解释这些话时说道：\Quote{值得注意的是，群众在选择这些人时没有分歧，宗徒们也没有被他们拒绝。但我们应知道他们有什么品级，领受了什么祝圣。是执事品吗？但这职务当时在教会中尚未存在。或者是司铎的职务？但那时还没有主教，只有宗徒，因此我认为清楚明显的是，那时还没有执事或司铎的名称。}\par

因此，我们也宣布，上述七位执事不应被理解为在奥迹中服务的执事，按照前面所阐述的教导，而是受委托为聚集者的共同利益管理分派事务的人，他们在这方面对我们而言也是一种慈善和热忱关怀需要者的典范。\par

% structure: p0046 source=h2 target=section reason=relative-heading-rank; base=h2

\section{第十七条}

既然不同教会的圣职人员离开了自己受祝圣的教会，投奔了其他主教，未经自己主教的同意便定居在其他教会中，从而表现出傲慢与不顺从；我们规定，自过去第四小纪的元月起，任何品级的圣职人员，未经自己主教的辞任书，无权被登记在另一教会的圣职名册中。今后凡不遵守此规定，而自取其辱并连累祝圣他的主教者，应与他本人及接纳他的主教一同撤职。\par

% structure: p0048 source=h2 target=section reason=relative-heading-rank; base=h2

\section{第十八条}

那些因蛮族入侵或因其他情况而出走的圣职人员，我们命令他们在原因消失后，或蛮族入侵结束后，返回自己的教会。不得无故长期离开。若有任何人未按本法规的规定返回，应将他切断共融，直到他返回自己的教会为止。收纳他的主教也受同样处罚。\par

% structure: p0050 source=h2 target=section reason=relative-heading-rank; base=h2

\section{第十九条}

监督教会的人，应当每日，尤其是在主日，教导全体圣职人员和百姓关于虔诚和正确宗教的言论，从圣经中收集关于真理的沉思和定论，不超越现在所定的界限，也不偏离承受天主的教父们的传统。若有关于圣经的任何争议，不得以不同于教会的光明导师们在其著作中阐释的方式解释，而要以此为荣，而不是凭自己的头脑编造，以免因缺乏技巧而偏离适宜之道。因为藉着上述教父们的教导，百姓认识何为善与可欲，以及何为无用与当弃绝，从而改善自己的生活，不被无知引导，而是专心致志于教义，注意不让恶事临到，并在对即将来临的惩罚的恐惧中成就自己的救恩。\par

% structure: p0052 source=h2 target=section reason=relative-heading-rank; base=h2

\section{第二十条}

主教不得在非所属的城市公开教导。若有被发现如此行，应终止其主教职务，但可履行司铎的职分。\par

% structure: p0054 source=h2 target=section reason=relative-heading-rank; base=h2

\section{第二十一条}

那些因违犯教规而犯罪，并因此应受完全且永久的罢黜者，被贬为平信徒。然而，若他们常将悔改置于眼前，自愿为其失宠并使自己与恩宠隔绝的罪而哀哭，他们仍可按照圣职人员的样式剪发。但若他们不愿服从本条教规，则必须如平信徒一样蓄发，因为他们是选择世俗共融而弃绝天上生活的人。\par

% structure: p0056 source=h2 target=section reason=relative-heading-rank; base=h2

\section{第二十二条}

凡因金钱而被祝圣者，无论主教或任何品级，而非经过考核和品行选择，我们命令将其罢黜，连同祝圣他们的人也一并罢黜。\par

% structure: p0058 source=h2 target=section reason=relative-heading-rank; base=h2

\section{第二十三条}

任何主教、司铎或执事，在授予无玷圣体时，不得向领受者索取任何费用。因为恩宠不可售卖，圣神的圣化也不得以金钱换取；而应将此恩赐简朴地分施给堪当领受的人。但若有圣职名册中的人向领圣体者提出要求，应将其罢黜，因其效仿了\Person{西满}的错误和邪恶。\par

% structure: p0060 source=h2 target=section reason=relative-heading-rank; base=h2

\section{第二十四条}

任何列入司铎名册的人或隐修士，不得参与赛马或出席戏剧表演。若有圣职人员被邀请参加婚宴，一旦游戏开始，应起身离开，因为这是我们教父们的教导。若有任何人被证实犯此过失，应停止或罢黜。\par

% structure: p0062 source=h2 target=section reason=relative-heading-rank; base=h2

\section{第二十五条}

此外，我们重申那条教规，命令\OriginalTerm{乡村}{\GreekText{ἀγροικικὰς}}堂区与\OriginalTerm{本地的}{\GreekText{ἐγχωρίους}}堂区应继续受原管辖的主教管辖；特别是若三十年连续管理而无异议。但若三十年内曾有或将有争议，受损害者可向省议会提出。\par

% structure: p0064 source=h2 target=section reason=relative-heading-rank; base=h2

\section{第二十六条}

若一位司铎因无知而缔结了非法婚姻，尽管他仍保留其席位权利，如神圣教规所规定，但他必须停止一切司铎职事。因为对这样的人赐予宽容已足够。他不适合祝福他人，因他自己需要照顾自己的伤口，因为祝福是圣化的灌注。但他怎能将自身所无的（因无知之罪）灌注给别人呢？因此，无论在公开还是私下，他都不能祝福，也不能分施基督的圣体，[也不能履行任何其他职务]；但应满足于荣耀的席位，向主哀哭，祈求其无知之罪得以赦免。因为显而易见，这邪恶的婚姻必须解除，他也不能与被剥夺司祭职执行的人有任何交往。\par

% structure: p0066 source=h2 target=section reason=relative-heading-rank; base=h2

\section{第二十七条}

凡列入圣职名册者，无论在城镇中居住或旅行时，不得穿着不合其身份的服装；而应穿着指定给圣职人员的服装。若有违反此条者，将被停职一周。\par

% structure: p0068 source=h2 target=section reason=relative-heading-rank; base=h2

\section{第二十八条}

既然我们了解到，在若干教会中，根据长期流传的习俗，葡萄被带到祭台上，司铎将其与无血祭献一同献祭，并同时分施给百姓，我们规定，今后任何司铎不得这样做，而只应将祭品单独分施给百姓，以振作其灵魂并赦免其罪过。至于作为初果献上的葡萄，司铎可单独祝福它们，并作为感恩之礼分施给求取的人，以感谢那按神圣旨意赐予果实、使我们身体增长和得到滋养的天主。若有圣职人员违反此规定，应被罢黜。\par

% structure: p0070 source=h2 target=section reason=relative-heading-rank; base=h2

\section{第二十九条}

\Place{迦太基}会议的一条教规说，祭台上的神圣奥迹只能由禁食者举行，唯一年中一日，即主设立晚餐的那日除外。当时，或许由于当地某些对教会有益的原因，甚至神圣教父们自己也使用了这一宽免。但既然没有什么能让我们放弃严格的遵守，我们规定应遵循宗徒和教父的传统；并决定在四旬期最后一周的星期四不可开斋，以免羞辱整个四旬期。\par

% structure: p0072 source=h2 target=section reason=relative-heading-rank; base=h2

\section{第三十条}

为谋教会之建立而愿行诸事，我等决意亦照料身处蛮族之司铎。故若彼等以为应超越宗徒法典关于不得假借虔诚与宗教借口离弃妻子之规定，而行超乎所命之事，因而与其妻合意禁绝同居，我等判定彼等绝不可再与妻同住，如此方可使彼等以完美之证示其许诺。然我等准此于彼等，唯因其禀性狭隘及异邦未定之习俗，别无他故。\par

% structure: p0074 source=h2 target=section reason=relative-heading-rank; base=h2

\section{第三十一则}

凡在宅内祈祷所举行神圣奥迹或施行圣洗之圣职人员，我等判定应得当地主教之同意。故若有圣职人员不如此行，则应予以撤职。\par

% structure: p0076 source=h2 target=section reason=relative-heading-rank; base=h2

\section{第三十二则}

既获悉在亚美尼亚地区，彼等仅于圣台献酒，行无血祭献者不掺水而为之，并以教会圣师若望·屈梭多模为其权威依据；屈梭多模在其对《玛窦福音》释义中云：\par

\Quote{祂复活后为何不饮水而以酒？是为拔除另一邪恶异端之根。因有某些人在奥迹中用水，以示祂既在传授奥迹时赐酒，又在复活后于无奥迹之简餐中仍用酒；祂说\InnerQuote{葡萄树的果实}，但葡萄树产酒，非水。}自此彼等以为此圣师否定圣祭中掺水之举。今为避免彼等此后在此点上陷于无知，我等揭示该圣父之正统见解。盖古有\OriginalTerm{水派}{Hydroparastatæ}之邪恶异端，其在祭献中以水代酒；此神圣者驳斥此种异端之可憎教导，指明其直接违背宗徒传统，故述上述之言。因彼曾命其受托牧职之教会，于无血祭献中掺水与酒，以彰显为普世生命及赎罪而从我等救主基督宝贵肋旁所流之水血混合；此外，凡灵光普照之教会，皆遵此神圣命令。\par

且我主基督按肉身之兄弟雅各，初受托耶路撒冷教会之座，以及凯撒勒雅教会总主教巴西略，其荣耀传遍普世，于书写交付我等关于奥迹祭献之指示时，明言神圣爵杯在神圣礼仪中应以水与酒祝圣。在迦太基集会之圣教父们亦明定：\Quote{在神圣奥迹中，除主之体血外，不得献其他，即按主亲自所定，乃掺水之酒与饼。}故若有主教或司铎不依宗徒所传而行此神圣事，且不以掺水之酒献祭，则应予以撤职，因其未完全彰显奥迹，且擅改所传之事。\par

% structure: p0080 source=h2 target=section reason=relative-heading-rank; base=h2

\section{第三十三则}

我等知在亚美尼亚地区，唯司祭后裔（仿效犹太人习俗）被委以圣职；甚至有些未削发者被派继承神圣法律之咏唱与诵读之职。故我等判定，自此以后，愿将任何人引入圣职者，不得顾及被祝圣者之出身；而应审查其是否（按圣规所定）堪当列入圣职名册，以便按教会升迁，无论其是否为司祭后裔。再者，不得允许任何人依圣职人员名册次序在读经台诵读，除非其已领受司祭削发及本牧者合法祝福。若有人被发现违反此等规定，则应予以断绝。\par

% structure: p0082 source=h2 target=section reason=relative-heading-rank; base=h2

\section{第三十四则}

此后，因司祭规条明言，结党或秘密结社之罪行既为世俗法律所禁，在教会中更应禁止；我等亦加意遵守：若发现任何圣职人员或隐修士结党或加入秘密结社，或图谋反抗主教或其他圣职人员，应一律剥夺其品秩。\par

% structure: p0084 source=h2 target=section reason=relative-heading-rank; base=h2

\section{第三十五则}

任何都主教在其教省主教逝世后，不得侵占或出售逝者之私产，或该教会之寡产；而应将其交予该主教区之圣职人员保管，直至选举出新主教，除非该教会已无圣职人员留存。若然，则都主教应保护该财产，不使之减少，待新主教任命时悉数交还。\par

% structure: p0086 source=h2 target=section reason=relative-heading-rank; base=h2

\section{第三十六则}

我等重申在天主护佑之帝国京城集会之一百五十位教父，与在加采东集会之六百三十位教父所颁之法规：判定君士坦丁堡之教座享有与古罗马教座同等之特权，于教会事务中应受同等尊崇，且位居其次。君士坦丁堡之后，依次为亚历山大里亚教座、安提约基雅教座，而后为耶路撒冷教座。\par

% structure: p0088 source=h2 target=section reason=relative-heading-rank; base=h2

\section{第三十七则}

鉴于在不同时期曾有外邦人的入侵，许多城市因此受制于不信者，致使一位城市主教在被祝圣后，可能无法就任其主教座、以司祭之秩序定居其中，并照常实行属于主教的祝圣及其他诸事。为保持对司祭职的尊敬与敬礼，且决不愿利用外邦人的侵害来破坏教会的权利，我们规定：凡如此被祝圣者，因上述原因未能定居于其主教座者，不受此事的任何损害，可保持其合法地位；他们可按教会法规祝圣不同等级的圣职人员，并在规定范围内使用其职务的权柄；并且，凡由他们出令之行政事项，均为有效且合法。因为当法律的精确遵守受到限制时，其职务的行使不应被必要之时所限制。\par

% structure: p0090 source=h2 target=section reason=relative-heading-rank; base=h2

\section{第三十八条}

我们也遵守教父们所立的规定，该规定如是决定：若任何城市因皇帝权柄而更新，或将要更新，则教会事务之秩序应遵循民事及公共模式。\par

% structure: p0092 source=h2 target=section reason=relative-heading-rank; base=h2

\section{第三十九条}

鉴于我们的弟兄与同工，塞浦路斯岛的主教若望，同其赫勒斯滂省的人民，因外邦人的侵袭，并为了免除异教徒的奴役，而只服从于最基督徒统治的权杖，已藉着仁慈天主的眷顾和我们的爱基督与虔诚的女皇的辛劳，从该岛迁出；我们决定：最初在厄弗所聚会的诸位神圣教父所赋予的特权，应不受任何革新而保留，即：新查士丁尼堡应享有君士坦丁堡的权利，凡被立为该地虔诚且最敬神主教者，应优先于赫勒斯滂省的所有主教，并按古老习惯由其属下主教选举。因为诸位神圣教父也曾努力维护各教会中既有的习惯，即基齐库斯城现任主教应服从于上述查士丁尼堡的教省总主教，以效法所有其他属于上述蒙天主喜爱的总主教若望的主教；按照习惯，基齐库斯城的主教亦应由若望主教祝圣。\par

% structure: p0094 source=h2 target=section reason=relative-heading-rank; base=h2

\section{第四十条}

因藉着远离生活的喧嚣与纷扰而依附于天主，实为非常有益；故我们不应未经审察就在适当时间之前接纳选择隐修生活者，而应遵守教父所传授的时限，以便我们能在达成明辨年龄后，以知识与判断的成果，将按照天主的生命之誓愿接纳为永远坚定。因此，凡将要接受隐修生活之轭者，年龄不应少于十岁；此事之审察取决于主教的决定，视其认为更长的时间是否更有利于此人进入并坚定于隐修生活。因为伟大的\Person{巴西略}在其神圣法规中规定，凡自愿奉献给天主并拥抱贞洁者，若已年满十七，应列入贞女之列；然而，我们遵循有关寡妇与女执事的先例，经过类推与比例考量，按上述时间接纳了选择隐修生活者。因为神圣宗徒记载：寡妇应在教会中被选举，年满六十岁者。但神圣教规规定，女执事应在四十岁时受祝圣，因为他们看到教会藉着天主的恩宠更为有力及强壮，并进一步前进；且他们看到信众在遵守天主诫命上的坚定与稳固。因此，我们也极其正确地理解此事，将恩宠的祝福赐予将要按天主进入此斗争者，如同迅速在他身上印上某种印记，从而引导他不至长久犹豫、退缩，甚至激发他选择并决志向善。\par

% structure: p0096 source=h2 target=section reason=relative-heading-rank; base=h2

\section{第四十一条}

凡在城镇或村庄中愿进入隐修院、并独自退隐照管自己的人，在进入隐修院修习独修生活之前，应怀着对天主的敬畏，在三年的时间内服从该房屋的长上，并合理地完成一切服从之事，以此显示他们对此生活的选择，并表明他们全心自愿地拥抱之；然后由该地方的长上（\OriginalTerm{院长}{\GreekText{προέδρος}}）审核；之后须在隐修院外再勇敢地忍受一年，以使他们的意向得以完全显明。因为藉此他们将充分而完全地证明，他们并非追求虚荣，而是因隐修生活本身的美善与尊贵而追求之。如此期限满后，若他们仍坚持同一意向而选择此生活，则应被封闭入内；此后非因共同利益或死亡逼近的其他迫切必要，且得该地主教的祝福，不得再随意走出此房屋。\par

凡未经上述原因而擅自走出其修院者，首先应强制将之禁闭于上述修院，然后以禁食及其他苦行治愈自身，因为他们应知道经上记载：\Quote{无人手扶犁而向后看的，适于天主的国。}\par

% structure: p0099 source=h2 target=section reason=relative-heading-rank; base=h2

\section{第四十二条}

那些被称为独修士、身穿黑袍、长发、游走于城市、与世俗男女交往、使其修道身份蒙受恶名者——我们规定：若他们愿意接受其他隐修士的会衣、剪短头发，则可被关闭在隐修院中，列于弟兄之列；但若他们不愿如此行，则应被逐出城市，被迫居住于旷野（\OriginalTerm{旷野}{\GreekText{ἐρήμους}}），他们的名称也正是由此而来。\par

% structure: p0101 source=h2 target=section reason=relative-heading-rank; base=h2

\section{第四十三条}

凡基督徒均可选择修会纪律的生活，抛却此生事务的烦扰波涛，进入隐修院，并按照隐修士的方式剃发，而不论他先前犯有何等过失。因为我们的救主天主说：\Quote{凡到我这里来的，我决不抛弃。}\par

因此，既然隐修的生活方式如同在石版上铭刻补赎的生活，凡真诚投奔者，我们都予以接纳；任何习俗均不得阻碍他实现其意愿。\par

% structure: p0104 source=h2 target=section reason=relative-heading-rank; base=h2

\section{第四十四条}

凡\Emph{隐修士}犯奸淫被定罪，或为婚姻共融及社交而娶妻者，应按教会法规受奸淫者之罚。\par

% structure: p0106 source=h2 target=section reason=relative-heading-rank; base=h2

\section{第四十五条}

我们得知，在某些女隐修院中，即将受神圣会衣者，先由带其前来者以丝绸及各色服饰、甚至黄金珠宝装扮，然后如此走近祭台，并在那里脱去这些财富的炫耀，随即领受会衣祝福，穿上黑袍；我们规定此后不得如此行。因为那已凭自己自由意志摒弃一切生活享乐、拥抱合乎天主的生活方式、并以坚固稳定的理由予以确认、因而来到隐修院者，不应再回忆那些早已遗忘的、属于这朽坏消逝之尘世的事物。因为这样做会在她们心中引起疑虑，使灵魂如汹涌波涛般烦乱，摇摆不定。此外，她们不应以哭泣来表现内心的沉重；但若几滴眼泪从眼中落下，正如很可能发生的那样，看见者可能以为这些眼泪是出于对苦修奋斗的\OriginalTerm{并非更}{\GreekText{μὴ} \GreekText{μᾶλλον}}情感（\OriginalTerm{情感}{\GreekText{διαθέσεως}}），而\OriginalTerm{比}{\GreekText{ἢ}}是因为她们离开尘世及世俗事物。\par

% structure: p0108 source=h2 target=section reason=relative-heading-rank; base=h2

\section{第四十六条}

选择苦修生活并定居于隐修院的女性，绝不可离开隐修院。但如果任何无法抗拒的必要迫使她们如此行，则需得到院长姆姆的祝福和许可；即使如此，她们也不可单独外出，而应与院中几位年长且卓越的隐修女一同外出，并听从院长命令。但她们绝不可留在外面。\par

同样，度隐修生活的男性，在紧急必要时，得凭受托管理者的祝福外出。\par

因此，凡违反我们现在所规定者，无论男女，均当受相称之罚。\par

% structure: p0112 source=h2 target=section reason=relative-heading-rank; base=h2

\section{第四十七条}

妇女不可在男隐修院中住宿，男子也不可在女隐修院中住宿。因为信友应当无可指摘、不引人跌倒，并应体面、正派、令天主喜悦地安排自己的生活。但若有人如此行，无论是圣职人员或平信徒，应予以绝罚。\par

% structure: p0114 source=h2 target=section reason=relative-heading-rank; base=h2

\section{第四十八条}

被擢升至主教尊位者的妻子，应经双方同意与丈夫分离；在他被祝圣为主教后，她应进入一所远离主教寓所的隐修院，并在那里享受主教供给。若被认为合适，她可被擢升至女执事尊位。\par

% structure: p0116 source=h2 target=section reason=relative-heading-rank; base=h2

\section{第四十九条}

我们重申神圣法规，规定凡经主教意愿祝圣的隐修院，应永久作为隐修院，其所属之物应归隐修院保留，不得再作为世俗住所，亦不得由任何人给予世俗人士。但若此前已有此类事发生，我们宣布其为无效；此后企图如此行者，应受教会法规之罚。\par

% structure: p0118 source=h2 target=section reason=relative-heading-rank; base=h2

\section{第五十条}

自此以后，任何人，无论是圣职人员或平信徒，均不得掷骰子赌博。若此后发现有人如此行，圣职人员应予以撤职，平信徒应予以绝罚。\par

% structure: p0120 source=h2 target=section reason=relative-heading-rank; base=h2

\section{第五十一条}

这神圣大公会议完全禁止那些所谓\Quote{戏子}及其\Quote{表演}，以及狩猎展览和舞台舞蹈。若有人轻视本法规，投身于任何禁止之事，圣职人员应予以撤职，平信徒应予以绝罚。\par

% structure: p0122 source=h2 target=section reason=relative-heading-rank; base=h2

\section{第五十二条}

在神圣四旬斋的所有日子，除安息日、主日及预报救主降生节外，应举行预祭圣体礼仪。\par

% structure: p0124 source=h2 target=section reason=relative-heading-rank; base=h2

\section{第五十三条}

鉴于属灵关系大于肉身亲缘；又因我们获悉，在某些地方，有人在救恩圣洗中成为孩童的代父母，随后与孩童的母亲（寡妇）缔结婚姻，我们规定今后不得如此行。但若有人在本法规之后仍行此事，他们必须首先中止此非法婚姻，然后受奸淫者之罚。\par

% structure: p0126 source=h2 target=section reason=relative-heading-rank; base=h2

\section{第五十四条}

圣经清楚地教导我们：\BibleQuote{不可亲近任何近亲，揭露他们的下体。} 背负天主者\Person{巴西略}在其教规中列举了一些被禁止的婚姻，并略去了大部分，而这两种方式都为我们提供了益处。因为他避免了许多不光彩的名称（免得用这些词玷污他的论述），用概括性的词语包含了不洁之事，以此大致向我们指明了被禁止的婚姻。但由于这种沉默以及难以理解哪些婚姻被禁止，问题变得模糊不清；我们认为更清楚地阐明这一点是好的，特此规定：从今以后，凡娶其父兄弟之女者；或父亲与儿子同娶母女；或父亲与儿子同娶两姐妹；或母亲与女儿同娶两兄弟；或两兄弟同娶两姐妹者，应受七年之教规处罚，除非他们公开脱离这种非法结合。\par

% structure: p0128 source=h2 target=section reason=relative-heading-rank; base=h2

\section{第五十五教规}

既然我们知道在\Place{罗马城}，在神圣的四旬期斋戒中，他们在星期六守斋，这违背了传承的教会惯例，神圣会议认为在\Place{罗马教会}中，那项\Quote{凡有圣职人员被发现于星期日或星期六守斋（除一次例外）者，应被撤职；若为平信徒，则应被绝罚}的教规也应坚定不移地遵守。\par

% structure: p0130 source=h2 target=section reason=relative-heading-rank; base=h2

\section{第五十六教规}

我们同样得知，在\Place{亚美尼亚}地区和其他地方，有些人在神圣四旬期的安息日和主日吃蛋和奶酪。因此，看来整个遍布天下的天主教会应遵循同一规则，完美地守斋，并且正如他们戒除一切宰杀之物一样，他们也应戒除蛋和奶酪，因为这些是那些我们戒除其肉的动物的产物。但如果有人不遵守此法，若为圣职人员，则撤职；若为平信徒，则绝罚。\par

% structure: p0132 source=h2 target=section reason=relative-heading-rank; base=h2

\section{第五十七教规}

不可将蜂蜜和牛奶献在祭台上。\par

% structure: p0134 source=h2 target=section reason=relative-heading-rank; base=h2

\section{第五十八教规}

凡属平信徒等级者，在有主教、司铎或执事在场时，不得自行分送神圣奥秘。若有人胆敢如此行，因违背既定之事，应被绝罚一周，从此让他学会不要自视过高。\par

% structure: p0136 source=h2 target=section reason=relative-heading-rank; base=h2

\section{第五十九教规}

圣洗圣事绝不可在房屋内的小堂举行；凡堪当领受无瑕光照者，应前往天主教会，在那里领受此恩赐。若有人被证违反我们所定者，若为圣职人员，撤职；若为平信徒，绝罚。\par

% structure: p0138 source=h2 target=section reason=relative-heading-rank; base=h2

\section{第六十教规}

既然\Person{宗徒}宣告，\BibleQuote{凡与主联合的，便是与主成为一灵}，显然，凡亲近主之仇敌的，便因与其亲近而成为一体。因此，那些假装被魔鬼附身，并因他们邪恶的举止而伪装显现其形状和外表的人，理当受到惩罚，并承受与那些真正被魔鬼附身者同样性质的痛苦和艰难，以使他们从魔鬼的（工作或更确切地说）势力中被解救出来。\par

% structure: p0140 source=h2 target=section reason=relative-heading-rank; base=h2

\section{第六十一教规}

那些投身于\OriginalTerm{占卜者}{soothsayers}、或被称为\OriginalTerm{百人队长者}{hecatontarchs}、或任何此类之人，以便从他们那里得知他们想要被揭示之事的人，凡此类人，按照教父们近来所颁布的法令，应受六年之教规处罚。那些为了娱乐和伤害纯朴之人而携带母熊或此类动物的人，也应受此惩罚；同样，那些算命、预言命运、族谱以及大量此类出自欺骗和欺诈之妄语的人，以及那些被称为\OriginalTerm{驱云者}{expellers of clouds}、\OriginalTerm{施咒者}{enchanters}、\OriginalTerm{护身符施与者}{amulet-givers}和\OriginalTerm{占卜者}{soothsayers}的人。\par

% structure: p0142 source=h2 target=section reason=relative-heading-rank; base=h2

\section{第六十二教规}

所谓的\OriginalTerm{朔日节}{Calends}，以及称之为\OriginalTerm{波塔节}{Bota}和\OriginalTerm{布鲁马利亚节}{Brumalia}，以及三月一日的群众集会，我们愿从信徒的生活中废除。还有妇女的公开舞蹈，这会带来许多伤害和恶果。此外，我们从基督徒的生活中驱逐那些以希腊人谎称为神祇的男女之名举行的舞蹈，这些舞蹈是按照古老而非基督徒的方式进行的；兹规定：从今以后，没有人可穿着女装，也没有妇女穿着适合男子的衣服。也不得戴喜剧、萨提尔或悲剧面具；也不得在压榨葡萄酒时呼求那可憎的\OriginalTerm{巴克科斯}{Bacchus}之名；也不得在将酒倒入坛中时（为了引人发笑），因无知和虚荣而做出源于疯狂欺骗之事。因此，那些在将来明知故犯这些所记载之事的人，若为圣职人员，我们命令撤职；若为平信徒，绝罚。\par

% structure: p0144 source=h2 target=section reason=relative-heading-rank; base=h2

\section{第六十三教规}

我们禁止在教会中公开诵读那些由真理的敌人伪造的殉道者列传，以侮辱基督的殉道者并引起听者不信；我们命令将这些书籍投入火中。但那些接受它们或视其为真的人，我们予以绝罚。\par

% structure: p0146 source=h2 target=section reason=relative-heading-rank; base=h2

\section{第六十四教规}

不允许平信徒公开辩论或教导，从而为自己僭取教导权柄，而应顺从主所设立的秩序，向那些领受教导恩宠的人侧耳，并受他们教导关于天主的事；因为在同一教会内，天主按宗徒的话设立了\Quote{不同的肢体}；而神学家额我略明智地解释这段经文时，称赞他们当时所遵循的秩序说：\Quote{弟兄们，我们尊重这秩序，我们持守这秩序。让这人是耳，那人是舌、手或任何别的肢体。让这人教导，让那人学习。}稍后又说：\Quote{顺服地学习，欣然充盈，敏捷地服务，我们不会都是那更活跃的肢体——舌，不会都是宗徒，不会都是先知，也不会都是讲解者。}又说：\Quote{你既是羊，为何自充牧人？你既是脚，为何要做头？你既列于士兵之数，为何企图做统帅？}在别处又说：\Quote{智慧吩咐：不要急于言语；不要与富人相比，你原是贫穷的；不要寻求比智者更聪明。}如果有人被发现削弱本规条，应予以绝罚四十天。\par

% structure: p0148 source=h2 target=section reason=relative-heading-rank; base=h2

\section{第65规条}

有人在月初于店铺和房屋前点火，并依照某种古代习俗愚昧癫狂地跳跃，我们命令从此停止这类行为。因此，凡做这事者，若为圣职人员，予以罢免；若为平信徒，予以绝罚。因为在《列王纪下》中记载：\BibleQuote{默纳舍在上主殿宇的两院中为上天的万军筑了祭坛，使他的儿子从火中经过，并施行占卜、巫术、飞鸟占卜，设立招魂者（或皮同）并增多了术士，行上主视为恶的事，惹祂发怒。}\par

% structure: p0150 source=h2 target=section reason=relative-heading-rank; base=h2

\section{第66规条}

从我们天主基督复活的神圣之日直到下一个主日，整整一周，信友在圣堂中应免于劳作，以圣咏、赞美诗和属灵歌曲在基督内喜乐，庆祝庆节，专心诵读圣经，并享受神圣奥迹；因为这样我们将与基督一同被高举，并与祂一同复活。因此，在上述日子里，不得举行任何赛马或任何公开表演。\par

% structure: p0152 source=h2 target=section reason=relative-heading-rank; base=h2

\section{第67规条}

圣经吩咐我们戒绝血、勒死之物和奸淫。因此，那些为了贪图口腹之欲，以任何技艺调制动物之血并吃食的人，我们予以适当惩罚。倘若今后有人胆敢以任何方式吃动物之血，若为圣职人员，予以罢免；若为平信徒，予以绝罚。\par

% structure: p0154 source=h2 target=section reason=relative-heading-rank; base=h2

\section{第68规条}

任何人不得毁坏或切割旧约或新约的书卷，或我们神圣且经认可的宣讲者与导师的著作，也不得将其交给书商或所谓的香料商，或交予任何其他类似的人销毁；除非该书卷已被书虫、水或其他方式损毁而无法使用。今后若有人被发现做这样的事，予以绝罚一年。同样，凡购买此类书卷者（除非他为自己保存，或为他人利益而保存）并试图毁坏它们，予以绝罚。\par

% structure: p0156 source=h2 target=section reason=relative-heading-rank; base=h2

\section{第69规条}

平信徒不得进入圣所（希腊文：至圣祭台），尽管按照某种古代传统，当皇帝想向造物主奉献礼物时，其权力和权威绝不被禁止进入。\par

% structure: p0158 source=h2 target=section reason=relative-heading-rank; base=h2

\section{第70规条}

妇女在举行神圣礼仪时不得发言；但按照宗徒保禄的话：\BibleQuote{她们应沉默。因为不准她们发言，而要服从，正如法律所说的。她们若愿意学习什么，可在家里问自己的丈夫。}\par

% structure: p0160 source=h2 target=section reason=relative-heading-rank; base=h2

\section{第71规条}

学习世俗法律的人不得采纳外邦人的习俗，也不得被诱导去戏院，也不得保持所谓的\Quote{\OriginalTerm{基莱斯特拉}{Cylestras}}，也不得穿着与一般习俗相悖的衣服；这适用于他们开始受训时、结业时以及整个学习期间。倘若从今以后有人胆敢违反本规条，予以绝罚。\par

% structure: p0162 source=h2 target=section reason=relative-heading-rank; base=h2

\section{第72规条}

正统信徒不可娶异端女子为妻，正统女子也不可嫁异端男子。但若发现有人已如此行，我们要求视该婚姻为无效，并予以解除。因为不宜将不应相混的混在一起，也不应将羊与狼结合，也不应将罪人的份与基督的份结合。若有人违反我们所颁布的，予以绝罚。但若有人至今仍为不信者，尚未列入正统羊群，他们之间已缔结合法婚姻，然后其中一方选择了正途，前来接受真理之光，而另一方仍被错误之链所束缚，不愿定睛注视神圣光芒，且不信的妻子愿意与信主的丈夫同居，或不信的丈夫愿意与信主的妻子同居，则不应分离，按神圣宗徒的话：\BibleQuote{因为不信的丈夫因妻子而成圣，不信的妻子因丈夫而成圣。}\par

% structure: p0164 source=h2 target=section reason=relative-heading-rank; base=h2

\section{第73规条}

既然赋予生命的十字架为我们显示了救恩，我们应当小心，藉以我们从古旧堕落中被拯救的那十字架应得到应有的尊荣。因此，我们在思想、言语、情感上向它表示敬礼（\OriginalTerm{敬礼}{\GreekText{προσκύνησιν}}），命令有些人放置在地面上的十字架图形务必完全移除，免得那为我们赢得的胜利标记被行走其上的人践踏而受亵渎。因此，我们规定，从今以后凡在路面上绘制十字架标记者，予以绝罚。\par

% structure: p0166 source=h2 target=section reason=relative-heading-rank; base=h2

\section{第74规条}

不允许在主的房屋或教堂内举行所谓的\OriginalTerm{爱宴}{Agapæ}，即爱宴，也不可在屋内进食，也不可铺床。若有人胆敢如此行，应停止，否则将被绝罚。\par

% structure: p0168 source=h2 target=section reason=relative-heading-rank; base=h2

\section{第七十五宪章}

我们规定，那些在教会中司歌唱之职者，不得使用无序的喊叫，也不可强逼本性呼喊，或采用任何不合宜且不适合教会的方式；而应以极大的专注和痛悔，向那位鉴察隐秘的天主献上圣咏。因为神圣的训谕曾教导以色列的子孙应虔诚。\par

% structure: p0170 source=h2 target=section reason=relative-heading-rank; base=h2

\section{第七十六宪章}

那些负责尊敬教会的人，不得在圣所界线内设立进食之处，也不可在那里提供食物，或进行其他买卖。因为我们的救主天主，当祂在肉身中支搭帐棚时，教导我们不可将祂父的家变成买卖之所。祂也倒出了兑换银钱者的小钱，并赶走了所有亵渎圣殿的人。因此，若有人犯下上述过错，应被绝罚。\par

% structure: p0172 source=h2 target=section reason=relative-heading-rank; base=h2

\section{第七十七宪章}

那些献身于宗教的人，无论是圣职人员还是隐修士，都不应与妇女一同在浴室中沐浴；任何基督徒男子或平信徒也不应如此行。因为这事受到外教人的严厉谴责。若有人在此事中被抓获，若他是圣职人员，应予以罢免；若他是平信徒，应予以绝罚。\par

% structure: p0174 source=h2 target=section reason=relative-heading-rank; base=h2

\section{第七十八宪章}

应当让那些受光照者背诵信经，并在一周的第五天向主教或司铎复述。\par

% structure: p0176 source=h2 target=section reason=relative-heading-rank; base=h2

\section{第七十九宪章}

既然我们宣认童贞女的天主诞生毫无产痛，因为它是无种子而发生的，并将此向整个羊群宣讲，那么我们就纠正那些因无知而做出与此不符之事的人。因此，鉴于有些人在我们天主基督的圣诞节次日，烹煮\OriginalTerm{细面粉}{\GreekText{σεμίδαλῖν}}，并互相分送，以敬礼无瑕童贞母性的\OriginalTerm{产子}{puerperia}为借口，我们规定，此后信徒不可再做此类事情。因为当我们用平常之事和与我们自身相似之事来界定和描述她无可言喻的产子时，这并非尊敬那超越思维和言语而在肉身中生育了不可理解之圣言的童贞女。因此，若有人此后被发现做此类事情，若他是圣职人员，应予以罢免；若他是平信徒，应予以绝罚。\par

% structure: p0178 source=h2 target=section reason=relative-heading-rank; base=h2

\section{第八十宪章}

若有主教、司铎、执事，或任何列入圣职名单者，或平信徒，没有极其严重的必要或困难的事务而长期不能到教会，但若身在本城，连续三个主日——三个星期——不去教会，若他是圣职人员，应予以罢免；若他是平信徒，应予以绝罚。\par

% structure: p0180 source=h2 target=section reason=relative-heading-rank; base=h2

\section{第八十一宪章}

我们听说，在有些地方，在三圣颂的赞美诗中，在\Quote{圣而不朽}之后加上了\Quote{为我们被钉十字架者，求祢垂怜我们}。由于这不符合虔诚，已被古代圣教父们从赞美诗中剔除，正如那些插入这些新词的暴力异端者也被逐出教会；我们同样确认我们圣教父们从前所虔诚建立的事，并绝罚那些在此法令之后，允许在教会中对此至圣赞美诗添加此项或其他任何内容的人；但若违反者属于司祭阶层，我们命令剥夺其司祭尊位；若他是平信徒或隐修士，应予以绝罚。\par

% structure: p0182 source=h2 target=section reason=relative-heading-rank; base=h2

\section{第八十二宪章}

在一些可敬圣像的画中，画有一只羔羊，前驱以手指向它，这被视为恩宠的预像，通过法律预先指示出我们真正的羔羊、我们的天主基督。因此，我们接受这些古代的预像和影子，作为真理的象征和赐予教会的范型，但更喜爱\BibleQuote{恩宠和真理}，视之为法律的满全。为使\BibleQuote{那圆满的}至少在色彩的描绘中显现在众人眼前，我们规定，从此以后，在图像中代之以古代羔羊的，是那除免世界之罪过的羔羊、我们的天主基督的人形像，使众人借此理解天主圣言降卑的深度，并使我们忆起祂在肉身中的交往、祂的苦难与救赎之死，以及祂为整个世界所完成的救赎。\par

% structure: p0184 source=h2 target=section reason=relative-heading-rank; base=h2

\section{第八十三宪章}

任何人不得将圣体给予死者的遗体；因为经上记载\BibleQuote{你们拿去吃}。但死者的遗体既不能\Quote{拿}，也不能\Quote{吃}。\par

% structure: p0186 source=h2 target=section reason=relative-heading-rank; base=h2

\section{第八十四宪章}

遵循教父们的教规，我们规定：对于婴儿，每当找不到可靠的证人肯定他们确已受洗，且因年龄幼小不能对所领受的入门奥迹给出满意答复时，他们应当毫无阻碍地受洗，以免这种疑虑剥夺他们获得这一净化之圣化的可能。\par

% structure: p0188 source=h2 target=section reason=relative-heading-rank; base=h2

\section{第八十五宪章}

我们从圣经中得知：\BibleQuote{凭两三个见证人的口，每句话都可成立。}因此，我们规定，凡由主人在三位证人面前释放的奴隶，应享有这份荣誉；因为那些当时在场的人将增强并稳固所给予的自由，并使所见证之事得以信守。\par

% structure: p0190 source=h2 target=section reason=relative-heading-rank; base=h2

\section{第八十六宪章}

那些为毁灭自己灵魂而供养并培育妓女的人，若他们是圣职人员，应予以绝罚并被罢免；若他们是平信徒，应予以绝罚。\par

% structure: p0192 source=h2 target=section reason=relative-heading-rank; base=h2

\section{第八十七宪章}

离弃自己丈夫的妇人，若另嫁他人，便是犯奸淫，按圣巴西略之教导，此乃从耶肋米亚先知处极佳地汇集而来：\Quote{若妇人成了别人的妻子，她的丈夫不可再回到她那里，但她既污秽了，必仍污秽。}又\Quote{奸夫的愚昧而不虔诚。}因此，若她似乎无故离开丈夫，丈夫应得宽恕，她应受惩罚。宽恕将赐予他，使他能与教会共融。但那离弃依法所娶的妻子，另娶他人者，按主的话，是犯奸淫。我们的教父们已规定，这类人须作\Quote{哭泣者}一年，\Quote{听者}二年，\Quote{俯伏者}三年，第七年与信友一同站立，并如此堪当领受祭献（若他们含泪做补赎）。\par

% structure: p0194 source=h2 target=section reason=relative-heading-rank; base=h2

\section{第八十八则}

任何人不得将牲畜带入圣堂，除非旅客在极端必要时，因无棚舍或歇息之处，而进入该圣堂。因为若未将牲畜带入，它将会死亡，而他因失去驮畜，无法继续旅程，将有死亡的危险。我们所受的教导是：安息日是为人设立的；因此，人的安全与舒适无论如何应置于首位。但若有人被查出非因上述必要而将牲畜带入圣堂，若是圣职人员，予以罢免；若是平信徒，予以绝罚。\par

% structure: p0196 source=h2 target=section reason=relative-heading-rank; base=h2

\section{第八十九则}

信友在救主苦难的日子，以斋戒、祈祷和心灵痛悔来度过，应当斋戒至大安息日的午夜；因为神圣的福音作者玛窦和路加已向我们表明，那（复活）是在深夜发生的，一位使用\OriginalTerm{安息日将尽}{\GreekText{ὀΨὲ} \GreekText{σαββάτων}}一词，另一位使用\OriginalTerm{清晨深时}{\GreekText{ὄρθρου} \GreekText{βαθέος}}一词。\par

% structure: p0198 source=h2 target=section reason=relative-heading-rank; base=h2

\section{第九十则}

我们从神圣的教父们领受了教规：为尊敬基督的复活，我们在主日不可跪拜。为避免我们忽略这规矩的圆满，我们向信友明确：在星期六晚祷时，司铎们进入祭台后（按惯例），直到主日傍晚，任何人不得跪着祈祷；在主日傍晚，在夜祷进入后，我们再屈膝向主献上祈祷。因为自安息日之后的夜晚——那是我主复活的前驱——我们开始以心神向天主唱赞美诗，引领我们的庆节从黑暗进入光明，因此在一整日和一整夜中，我们庆祝复活。\par

% structure: p0200 source=h2 target=section reason=relative-heading-rank; base=h2

\section{第九十一则}

凡给堕胎药物，以及接受毒药杀死胎儿者，处以杀人罪的惩罚。\par

% structure: p0202 source=h2 target=section reason=relative-heading-rank; base=h2

\section{第九十二则}

神圣的教务会议规定：那些以婚姻为名抢夺妇女的人，以及以任何方式协助抢婚者，若是圣职人员，即丧失其品级；若是平信徒，则予以绝罚。\par

% structure: p0204 source=h2 target=section reason=relative-heading-rank; base=h2

\section{第九十三则}

若一个男人的妻子离家出走、下落不明，她在未确知丈夫死亡之前与另一人同居，她便是个淫妇。军人的妻子，若在丈夫下落不明时另嫁，也是如此；以及那些因丈夫流离而不再等待他们回来的妻子，也是如此。但在此情形中有些可原谅之处，因对其死亡的怀疑变得非常大。但若她无知地嫁给一个当时被其妻子离弃的男人，后来因他的原配返回而被休弃，她的确犯了淫乱，但出于无知；因此她不禁止再婚，但最好保持现状。若一个军人在长时间后返回，发现他的妻子因他长期不在而另嫁他人，他若愿意，可以接回自己的妻子，并因无知而宽恕她和那与她进行第二次婚姻的男人。\par

% structure: p0206 source=h2 target=section reason=relative-heading-rank; base=h2

\section{第九十四则}

教规对那些起异教徒誓言的人处以惩罚，我们规定对他们予以绝罚。\par

% structure: p0208 source=h2 target=section reason=relative-heading-rank; base=h2

\section{第九十五则}

那些从异端归向正统信仰、并加入应得救者之列的人，我们按下列规则与惯例接纳他们。亚略派、马其顿派、诺瓦替安派（他们自称洁净派）、左派、十四日派（或称四日派）以及阿波利拿里派，我们是在他们呈交证书并诅咒每一个不与神圣天主宗徒教会持相同信仰的异端后接纳他们：然后首先用圣油敷抹他们的额、眼、鼻、口和耳；当我们封印时，说：\Quote{圣神恩赐的印记。}\par

但关于保禄派，天主教会已规定他们必须重新受洗。欧诺米派（他们以一次浸入施行洗礼）、孟他努派（此处称为弗里吉亚派）、撒伯流派（他们认为圣子与圣父相同，并在其他严重问题上犯错），以及所有其他异端——因为这里有许多异端，尤其来自迦拉达地区的——所有这些愿意归向正统信仰的人，我们按外邦人接纳他们。第一天，我们使他们成为基督徒；第二天，成为慕道者；第三天，我们为他们驱魔，同时三次向他们面孔和耳朵吹气；然后我们收录他们，让他们在圣堂内停留并聆听圣经；之后我们为他们施洗。\par

至于摩尼派、瓦伦提尼派、马西昂派以及所有类似异端的人，必须呈交证书，诅咒各自的异端，以及聂斯多略、欧提克斯、狄奥斯库若、塞维鲁和这类异端的其他首领及其同党，以及所有前述异端；如此，他们才能分享圣体圣事。\par

% structure: p0212 source=h2 target=section reason=relative-heading-rank; base=h2

\section{第九十六则}

凡藉圣洗穿上基督的人，已宣告要效法祂在肉身所行的生活方式。因此，凡以损害观者目光的方式装饰和梳理头发，即用狡诈编结的辫子，从而以此方式为不稳定的灵魂设下圈套者，我们以父亲般的照管，藉适当的惩罚予以医治：训练并教导他们度节制的生活，为使他们弃绝物质事物的欺骗与虚荣，持续将心思投注于有福且无害的生活，并在敬畏中持守纯结、圣善的交谈，从而藉生活的纯洁尽可能接近天主；并修饰内在的人而非外表，以德行及善良无瑕的品行，使自己在敌人左手的遗毒中不留丝毫残余。若有人违反本教规，应予以绝罚。\par

% structure: p0214 source=h2 target=section reason=relative-heading-rank; base=h2

\section{第九十七则}

凡与妻子行房或以任何其他方式不尊重圣地，使之沦为世俗，并轻蔑对待而仍停留在其中者，我们命令将所有这样的人都逐出，甚至逐出位于可敬圣堂内的望教者住所。若有人不遵守这些指示，他若是圣职人员，应予以免职；若是平信徒，应予以绝罚。\par

% structure: p0216 source=h2 target=section reason=relative-heading-rank; base=h2

\section{第九十八则}

凡将已许配给另一尚在世的男子的女子引入婚姻结合者，应承担奸淫之责。\par

% structure: p0218 source=h2 target=section reason=relative-heading-rank; base=h2

\section{第九十九则}

我们还进一步得知，在\Place{亚美尼亚}地区，某些人在圣所内煮肉块，并将部分分给司铎，按犹太人的方式分发。因此，为保持圣堂不受玷污，我们规定：任何司铎不得从献祭者手中强取分割的肉块；他们应满足于献祭者乐意给予的部分；此外，我们还规定此类献祭应在圣堂外进行。若有人不如此行，应予以绝罚。\par

% structure: p0220 source=h2 target=section reason=relative-heading-rank; base=h2

\section{第一百则}

智慧吩咐说：\Quote{你的眼目要注视正直的事，} \Quote{并要尽心保守你的心。}因为肉身的感官容易将其印象带入灵魂。因此，我们命令：今后绝不得制作任何图画，无论是绘画还是其他方式，凡吸引眼目、败坏心思、并刺激人燃起卑劣欲乐者。若有人企图这样做，应予以绝罚。\par

% structure: p0222 source=h2 target=section reason=relative-heading-rank; base=h2

\section{第一百零一则}

伟大而神圣的宗徒保禄高声宣告：按天主的肖像受造的人，乃是基督的身体和圣殿。因此，他超越一切有灵之物，藉救赎的苦难达到属天的尊位，吃喝基督，便在各方面适于永生，藉参与天主的恩宠圣化自己的灵魂和身体。因此，若有人希望在聚会时分享无玷的圣体，并奉献自己为领圣体，应走近前来，将双手交叉成十字架的形状，如此领受恩宠的共融。但那些不用手，而是用金或其他材料制作的器皿来承接神圣恩赐，并以此领取无玷共融的人，我们绝不允许他们前来，因为他们更喜爱无生命和低等的物质，而非天主的肖像。若有人被发现向携带此类器皿的人分授无玷共融，则他与携带者一同受绝罚。\par

% structure: p0224 source=h2 target=section reason=relative-heading-rank; base=h2

\section{第一百零二则}

凡从天主领受释放与束缚权柄者，理应考量罪的严重性与罪人的悔改准备，并对症下药，以免因对每一点判断不当而未能治愈病者。因为罪的疾病并非单一，而是各种多样，滋生许多有害的枝节，由此散布许多邪恶，且持续蔓延，直至被医师的能力所遏制。因此，声称精通属灵医术的人，首先应考虑犯罪者的心态，看他是否趋向健康，抑或相反，因自己的行为招致疾病，并要观察他如何在期间照料自己的生活方式。若他未抗拒医师，且因施涂的药物而使灵魂的溃疡加重，则应按他所堪当的给予仁慈。因为全部账目乃在天主与受托牧职者之间，为领回迷失的羊，并治愈被蛇所伤者；既不使他们坠入绝望的深渊，也不松开缰绳，使之放纵或轻视生命；而是以某种方式，或严厉收敛，或更温和软性的药物，抵挡此疾病，竭力治愈溃疡，同时检查他的悔改果实，并明智地管理那被召向更高光照的人。因为我们应当知道两件事，即属于严格的事与属于惯例的事，并在那些不适合更高事物的人身上，遵循传统的形式，正如圣巴西略所教导我们。\par

\SourceRecord{Translated by Henry Percival.；Nicene and Post-Nicene Fathers, Second Series；Vol. 14.；Edited by Philip Schaff and Henry Wace.；Buffalo, NY: Christian Literature Publishing Co.,；1900.；<http://www.newadvent.org/fathers/3814.htm>.}
